\documentclass[11pt,a4paper]{amsart}
\usepackage{amsmath,amssymb,amsthm}
\usepackage[margin=2.6cm]{geometry}
\usepackage{booktabs}
\usepackage[colorlinks=true,linkcolor=blue,citecolor=blue,urlcolor=blue]{hyperref}

\newtheorem{theorem}{Theorem}[section]
\newtheorem{proposition}[theorem]{Proposition}
\newtheorem{lemma}[theorem]{Lemma}
\newtheorem{corollary}[theorem]{Corollary}
\theoremstyle{definition}
\newtheorem{conjecture}[theorem]{Conjecture}
\newtheorem{remark}[theorem]{Remark}
\newtheorem{example}[theorem]{Example}

\newcommand{\Gg}{\mathcal{G}}
\newcommand{\Gs}{\mathcal{G}_{\star}}
\newcommand{\R}{\mathbb{R}}
\newcommand{\eps}{\varepsilon}
\DeclareMathOperator{\arsh}{artanh}

\title[Sharp remainders and optimal localization for Gurland's ratio]{Sharp remainder
estimates and optimal localization for Gurland's ratio of the gamma function}

\author{}
\date{\today}

\subjclass[2020]{Primary 33B15; Secondary 11M35, 26D07, 41A58}
\keywords{Gurland's ratio, gamma function, Hurwitz zeta function, Jensen gap,
complete monotonicity, remainder estimates}

\begin{document}

\begin{abstract}
Gurland's ratio $\Gg(x,y)=\Gamma(x)\Gamma(y)/\Gamma^{2}\big(\tfrac{x+y}{2}\big)$ is the
exponential of the Jensen gap of $\log\Gamma$ at the pair $(x,y)$. We show that this gap
admits an exact power-series expansion in powers of $\big(\tfrac{x-y}{2}\big)^{2}$ whose
coefficients are Hurwitz zeta values --- a classical identity, recorded in this form by
Srivastava --- and that the expansion converges \emph{for every} $x,y>0$: the natural
radius of convergence is the distance from the midpoint
$\tfrac{x+y}{2}$ to the pole of $\Gamma$ at the origin, so the region of convergence is
exactly the open quadrant. This answers in the affirmative two open problems recently posed
by Wi\'sniewska (arXiv:2512.07028) and shows that the sufficient condition $Q<1$ she imposed
is redundant. We complement the expansion with two-sided remainder estimates whose
geometric rate $\rho=|x-y|/(2+x+y)$ we prove optimal; we show that the previously used rate
$Q$ satisfies $\rho\le Q$ with equality exactly when $x=y$, and that $Q$ is unbounded. We
give the complete asymptotics of the ratio for large arguments in terms of the binary
entropy function, extend the expansion to any number of variables, and derive applications
to the beta function and to certified computation. Our main result localizes the
mean-value parameter $t(x,y)$ of Wi\'sniewska; the two-sided endpoint bound that defines
$t$ appears already in the work of Yang and Zheng. We determine the asymptotics of $t$ up
to order $(x-y)^{4}$ and prove the sharp bound
\[
  \frac{t(x,y)-\sqrt{xy}}{\frac{x+y}{2}-\sqrt{xy}}\;>\;\frac{2}{3},
\]
the constant $\tfrac23$ being optimal and not attained. The proof is self-contained and
uses neither of the two sufficient conditions (A) and (B) under which the weaker midpoint
property $t>\tfrac12\big(\sqrt{xy}+\tfrac{x+y}{2}\big)$ had been established; we also
record an independent, elementary route that yields the bound on an explicit range. We
conjecture that conditions (A) and (B) cover the whole quadrant. All results are
accompanied by high-precision numerical verification.
\end{abstract}

\maketitle

\section{Introduction}

\subsection{Background}
For $x,y>0$ Gurland's ratio \cite{Gurland56} is
\begin{equation}\label{eq:G}
  \Gg(x,y)=\frac{\Gamma(x)\,\Gamma(y)}{\Gamma^{2}\big(\frac{x+y}{2}\big)},
\end{equation}
and Gurland's inequality states that $\Gg(x,y)>1$ for $x\ne y$. Since $\log\Gamma$ is
strictly convex on $(0,\infty)$ (Bohr--Mollerup), $\log\Gg(x,y)$ is exactly the
\emph{Jensen gap}
\begin{equation}\label{eq:jensen}
  \log\Gg(x,y)=\log\Gamma(x)+\log\Gamma(y)-2\log\Gamma\Big(\frac{x+y}{2}\Big)\;\ge\;0,
\end{equation}
so Gurland's ratio is the exponential of the second-order Jensen gap of $\log\Gamma$.
Following \cite{Wisniewska25} we also consider the modified ratio
\begin{equation}\label{eq:Gs}
  \Gs(x,y)=\frac{\Gamma(1+x)\,\Gamma(1+y)}{\Gamma^{2}\big(1+\frac{x+y}{2}\big)}
            =\frac{4xy}{(x+y)^{2}}\,\Gg(x,y).
\end{equation}
The ratio \eqref{eq:G} and its relatives were studied by Gautschi \cite{Gautschi74},
Kairies \cite{Kairies83}, Merkle \cite{Merkle05}, Alzer \cite{Alzer01} and Chen--Choi
\cite{ChenChoi17}; see \cite{LaforgiaNatalini13} for a survey.

Writing
\begin{equation}\label{eq:uv}
  u:=\frac{x-y}{2},\qquad v:=\frac{x+y}{2},\qquad\text{so that }x=v+u,\quad y=v-u,
\end{equation}
Wi\'sniewska \cite[Thm.~1]{Wisniewska25}, starting from the Weierstrass product, proved the
finite expansion
\begin{equation}\label{eq:Wthm}
  \log\Gs(x,y)=\sum_{k=1}^{m-1}\frac{1}{k}\Big(\frac{x-y}{2}\Big)^{2k}
     \zeta\Big(2k,1+\frac{x+y}{2}\Big)+R^{m}_{n}(x,y),\qquad m\ge2,
\end{equation}
with the remainder estimate
\begin{equation}\label{eq:Wrem}
  0\le R^{m}_{n}(x,y)\le \eps^{m}_{n}(x,y)
  =\frac1m\Big(\frac{x-y}{2}\Big)^{2m}\sum_{n=1}^{\infty}\frac{1}{(n+x)^{m}(n+y)^{m}},
\end{equation}
where $\zeta(s,a)=\sum_{n\ge0}(n+a)^{-s}$ is the Hurwitz zeta function. From
\eqref{eq:Wthm}--\eqref{eq:Wrem} she deduced that the condition
\begin{equation}\label{eq:Q}
  Q(x,y):=\frac{|x-y|}{2\,(1+\sqrt{xy})}<1
\end{equation}
implies $\log\Gs(x,y)=S_{\infty}(x,y)$, where $S_\infty=\lim_{m\to\infty}S_m$ and $S_m$
denotes the partial sum in \eqref{eq:Wthm} \cite[Cor.~1]{Wisniewska25}; the bilateral
inequality
\begin{equation}\label{eq:Wbil}
  \Big(\frac{x-y}{2}\Big)^{2}\zeta\Big(2,1+\frac{x+y}{2}\Big)\le\log\Gs(x,y)
  \le\Big(\frac{x-y}{2}\Big)^{2}\zeta\big(2,1+\sqrt{xy}\,\big)
\end{equation}
\cite[Cor.~2]{Wisniewska25}. The two-sided bound \eqref{eq:Wbil} is, however, not new.
Indeed, with the base point $1$, $p=x$ and $q=y$, Corollary~11 and Remark~5 of Yang and
Zheng \cite{YangZheng19} state that
\begin{equation}\label{eq:YZ}
  \exp\Big[\frac{(p-q)^{2}}{4}\psi'(1+r_{1})\Big]
  \;\le\;\frac{\Gamma(1+p)\,\Gamma(1+q)}{\Gamma^{2}\big(1+\frac{p+q}{2}\big)}
  \;\le\;\exp\Big[\frac{(p-q)^{2}}{4}\psi'(1+r_{2})\Big],
\end{equation}
with the best constants $r_{1}=\frac{p+q}{2}$ and $r_{2}=\sqrt{pq}$ whenever $p,q>0$; as
$\psi'=\zeta(2,\cdot)$ and $\big(\frac{x-y}{2}\big)^{2}=u^{2}$, this is exactly
\eqref{eq:Wbil}. Since $\zeta(2,\cdot)$ is continuous and strictly decreasing,
\eqref{eq:Wbil} yields at once a unique parameter
$t=t(x,y)\in\big(\sqrt{xy},\tfrac{x+y}{2}\big)$ with
\begin{equation}\label{eq:Wt}
  \log\Gs(x,y)=\Big(\frac{x-y}{2}\Big)^{2}\zeta(2,1+t)
  \qquad\text{\cite[Cor.~3]{Wisniewska25}.}
\end{equation}
Thus Corollaries 2 and 3 of \cite{Wisniewska25} follow from \cite{YangZheng19}. The
question that remains is quantitative --- \emph{where}, inside
$\big(\sqrt{xy},\frac{x+y}{2}\big)$, the point $t$ lies --- and it is this question that
Theorems \ref{thm:loc} and \ref{conj:sharp} answer.
She then posed \cite[\S2]{Wisniewska25} three problems.

\smallskip
\noindent\textbf{Open problem 1.} Does $R^{m}_{n}(x,y)\to0$ hold for \emph{every} $x,y>0$?

\noindent\textbf{Open problem 2.} Does $\log\Gs(x,y)=S_{\infty}(x,y)$ hold in general, and
if not, how large can the difference be?

\noindent\textbf{Open problem 3.} How is $t(x,y)$ located in the interval
$\big(\sqrt{xy},\tfrac{x+y}{2}\big)$? What is its distance to the midpoint of that interval?
\normalsize

\subsection{Relation to the literature, and what is new here}
The identity in Theorem \ref{thm:main} is a symmetrised Taylor expansion and is
\emph{classical}. Since $\psi^{(n)}(z)=(-1)^{n+1}n!\,\zeta(n+1,z)$ for
$\operatorname{Re}z>0$ \cite[\S5.15]{DLMF}, the Taylor series of $\log\Gamma$ about $a>0$ is
\[
  \log\Gamma(a+w)=\log\Gamma(a)+\sum_{n\ge1}\frac{(-1)^{n}}{n}\zeta(n,a)\,w^{n},
  \qquad |w|<a ,
\]
whose even part is exactly \eqref{eq:main1}; equivalently,
\begin{equation}\label{eq:classical}
  \log\Gamma(a+w)+\log\Gamma(a-w)-2\log\Gamma(a)
   =\sum_{k\ge1}\frac{\zeta(2k,a)}{k}\,w^{2k},\qquad |w|<a ,
\end{equation}
and it is in this form, as a sum of a series of Hurwitz zeta values, that the identity is
recorded in the literature (see, e.g., \cite{Srivastava88}) and in monographs on zeta
series \cite{SrivastavaChoi01}; the classical relation between the gamma function and the
Hurwitz zeta function goes back at least to \cite{Berndt85}. We therefore claim no novelty
for \eqref{eq:classical}, nor for its specialization $a=1+v$, $w=u$, which is
\eqref{eq:main2}. What is new here is
(i) the \emph{sharp remainder theory} of the resulting convergent expansion: the two-sided
estimate \eqref{eq:rembnd} with the \emph{optimal} geometric rate $\rho$, together with the
comparison with the previously used rate $Q$; (ii) the explicit observation that this
specialization makes Wi\'sniewska's expansion \eqref{eq:Wthm} converge on the whole
quadrant and her hypothesis $Q<1$ redundant, which settles Open Problems 1 and 2 of
\cite{Wisniewska25}; (iii) the quantitative localization theory of the trigamma mean $t$:
the bracket \eqref{eq:locbnd}, the expansion \eqref{eq:locasym} as $c\to0$, and the sharp
global bound $\theta>\frac23$ of Theorem \ref{conj:sharp}, whose best constant $\frac23$ is
not a corollary of the estimates of Theorem \ref{thm:loc} but follows from a separate
global argument in four steps; (iv) the large-argument asymptotics of
Theorem \ref{thm:asym}; (v) the multivariate extension of Theorem \ref{thm:multi}; and
(vi) the applications of Section \ref{sec:app}.

Gurland's ratio has also been studied intensively by asymptotic methods. Merkle
\cite{Merkle05} derived bounds and asymptotic expansions; Alzer \cite{Alzer97} and Qi
\cite{Qi10} gave further bounds and a survey; Tian and Yang \cite{TianYang21} established
the asymptotic expansion
\[
  \frac{\Gamma(x+p)\,\Gamma(x+q)}{\Gamma\big(x+\frac{p+q}{2}\big)^{2}}
   =\exp\Big[\sum_{k=1}^{n}\frac{B_{2k}(s)-B_{2k}\big(\frac12\big)}
   {k(2k-1)\,(x+r_{0})^{2k-1}}+R_{n}(x;p,q)\Big],\qquad x\to\infty ,
\]
where $p,q\in\R$, $w=|p-q|\ne0$, $s=\frac{1-w}{2}$, $r_{0}=\frac{p+q-1}{2}$ and $B_{2k}$
are the Bernoulli polynomials, and proved that $(-1)^{n}R_{n}$ is completely monotone when
$w<1$, which yields sharp bounds; the remainder was analysed further by Yang and Tian
\cite{YangTian24} for the generalized ratio
$\frac{\Gamma(x+a)\Gamma(x+b)}{\Gamma(x+c)\Gamma(x+d)}$ with $a+b=c+d$. These expansions
are \emph{asymptotic in $1/x$} with Bernoulli-polynomial coefficients, whereas
\eqref{eq:main1}--\eqref{eq:main2} are \emph{exact and convergent} in powers of
$\big(\frac{x-y}{2}\big)^{2}$ with Hurwitz-zeta coefficients. The two descriptions are
complementary: the former is adapted to large arguments, the latter to arguments that are
close together, and only the latter can be summed to the exact value. Refinements of
Gurland's formula for $\pi$ were given by Lin \cite{Lin13}.

Two further strands of the literature are close to the localization problem studied here.
Yang and Zheng \cite{YangZheng19} gave necessary and sufficient conditions for four classes
of ratios of gamma functions to be logarithmically completely monotonic; their
Corollary~11 and Remark~5 yield the two-sided bound \eqref{eq:YZ}, and their
Propositions~3 and~4 compare double integrals of $\psi^{(n)}$ over rectangles, which is the
same weighted-average picture that underlies \eqref{eq:locbnd}. In \cite{YangZheng16} the
mean $\psi_{n}^{-1}\big(\frac{1}{b-a}\int_{a}^{b}\psi_{n}(x+t)\,dt\big)-x$, generated by a
\emph{uniform} average of a polygamma function, is shown to increase from $\min(a,b)$ onto
$\frac{a+b}{2}$; the parameter $t$ of \eqref{eq:Wt} is the corresponding mean for a
\emph{triangular} average of $\psi'$ (see Remark \ref{rem:equiv}), so the two objects are of
the same type but not the same, and neither the existence of an interior sharp constant nor
its value follows from \cite{YangZheng16}. Very recently, Moustafa and Mahmoud
\cite{MoustafaMahmoud26} obtained bounds for Gurland-type ratios in terms of the generalized
Nielsen beta function. The sharp remainder theory of \eqref{eq:main1} and the quantitative
localization theory developed below do not appear in these works.

\subsection{Results of the present paper}
Our starting point is that \eqref{eq:jensen} is a second-order symmetric difference of
$\log\Gamma$, so the subject is governed by the Taylor expansion of $\log\Gamma$ about the
midpoint $v$. Since the singularities of $\log\Gamma$ nearest to $v>0$ are the poles of
$\Gamma$ at $0,-1,-2,\dots$, the Taylor series about $v$ converges on the disc $|z-v|<v$,
and $x=v+u$, $y=v-u$ both lie in that disc precisely when $|u|<v$, i.e.\ precisely when
$x,y>0$. Thus the expansion converges on the whole quadrant and no additional hypothesis is
required.

\begin{theorem}[Exact expansion]\label{thm:main}
For all $x,y>0$,
\begin{align}
  \log\Gg(x,y)&=\sum_{k=1}^{\infty}\frac{1}{k}\Big(\frac{x-y}{2}\Big)^{2k}
      \zeta\Big(2k,\frac{x+y}{2}\Big),\label{eq:main1}\\[1mm]
  \log\Gs(x,y)&=\sum_{k=1}^{\infty}\frac{1}{k}\Big(\frac{x-y}{2}\Big)^{2k}
      \zeta\Big(2k,1+\frac{x+y}{2}\Big).\label{eq:main2}
\end{align}
Both series converge absolutely. Moreover \eqref{eq:main1} converges for exactly those
$(x,y)\in\R^{2}$ with $|x-y|<x+y$, and \eqref{eq:main2} for exactly those with
$|x-y|<2+x+y$; on the quadrant $x,y>0$ both hold identically, and the two regions are
optimal.
\end{theorem}

\begin{corollary}\label{cor:open12}
For every $x,y>0$ and every $m\ge1$,
\begin{align*}
  R^{m}_{n}(x,y)&=\log\Gs(x,y)-\sum_{k=1}^{m-1}\frac{1}{k}\Big(\frac{x-y}{2}\Big)^{2k}
     \zeta\Big(2k,1+\frac{x+y}{2}\Big)\\
  &=\sum_{k=m}^{\infty}\frac{1}{k}\Big(\frac{x-y}{2}\Big)^{2k}
     \zeta\Big(2k,1+\frac{x+y}{2}\Big)\;\ge\;0 .
\end{align*}
Consequently $R^{m}_{n}(x,y)\downarrow0$ and $\log\Gs(x,y)=S_{\infty}(x,y)$ for
\emph{every} $x,y>0$: Open Problems 1 and 2 both have affirmative answers. In particular
the hypothesis $Q(x,y)<1$ in \cite[Cor.~1]{Wisniewska25} is redundant.
\end{corollary}

The next result quantifies the convergence; the relevant rate is not \eqref{eq:Q} but
\begin{equation}\label{eq:rho}
  \rho(x,y):=\frac{|x-y|}{2+x+y}\quad(\Gs),\qquad
  \rho_{\Gg}(x,y):=\frac{|x-y|}{x+y}\quad(\Gg).
\end{equation}

\begin{theorem}[Sharp remainder estimates]\label{thm:rem}
Let $x,y>0$, $m\ge1$, $u=\tfrac{x-y}{2}$, $v=\tfrac{x+y}{2}$,
$\rho=\rho(x,y)=|x-y|/(2+x+y)$ and
$R_m:=\log\Gs(x,y)-\sum_{k=1}^{m-1}\frac{u^{2k}}{k}\zeta(2k,1+v)$. Then
\begin{equation}\label{eq:rembnd}
  \frac{u^{2m}}{m}\,\zeta(2m,1+v)\;\le\;R_m\;\le\;
  \frac{\rho^{2m}}{m\,(1-\rho^{2})}\Big(1+\frac{1+v}{2m-1}\Big).
\end{equation}
Moreover, if $x\ne y$, then $\lim_{m\to\infty}\frac{1}{m}\log R_m=2\log\rho$, so the rate
$\rho$ in \eqref{eq:rembnd} is optimal. The same holds for $\Gg$ with $\rho$ replaced by
$\rho_{\Gg}=|x-y|/(x+y)$ and $\zeta(2m,1+v)$ by $\zeta(2m,v)$.
\end{theorem}

\begin{proposition}\label{prop:compare}
For all $x,y>0$ one has $\rho(x,y)\le Q(x,y)$, with equality iff $x=y$. Moreover
$\rho(x,y)<1$ for all $x,y>0$, whereas $\sup_{x,y>0}Q(x,y)=+\infty$; for instance
$Q(x,1)=\frac{x-1}{2(1+\sqrt x)}\sim\frac{\sqrt x}{2}$ as $x\to\infty$. Hence $\{Q<1\}$ is a
proper subset of the quadrant, and Corollary~1 of \cite{Wisniewska25} does not cover the
whole range of validity of the expansion.
\end{proposition}

\begin{corollary}\label{cor:bil}
For all $x,y>0$ and all $m\ge1$,
\[
  S_m(x,y)\;\le\;\log\Gs(x,y)\;\le\;S_m(x,y)+
  \frac{\rho^{2m}}{m(1-\rho^{2})}\Big(1+\frac{1+v}{2m-1}\Big),
\]
where $S_m(x,y)=\sum_{k=1}^{m-1}\frac{1}{k}\big(\frac{x-y}{2}\big)^{2k}\zeta(2k,1+v)$.
Taking $m=2$ improves the upper bound of \eqref{eq:Wbil} whenever
$\rho<\frac{\sqrt5}{3}$, while $S_m$ increases to $\log\Gs(x,y)$ as $m\to\infty$.
\end{corollary}

We also determine the behaviour of the ratio at large arguments. Recall
\begin{equation}\label{eq:Lambda}
  \Lambda(\rho):=(1+\rho)\log(1+\rho)+(1-\rho)\log(1-\rho),\qquad |\rho|<1 .
\end{equation}
One has $\Lambda(\rho)=\sum_{k\ge1}\rho^{2k}/\big(k(2k-1)\big)$, $\Lambda(\rho)>0$ for
$\rho\ne0$, and $\Lambda(\rho)=2\rho\arsh\rho+\log(1-\rho^{2})$. As $\Lambda$ vanishes to
second order at the origin it is convenient to divide it by $\rho^{2}$: we set
\begin{equation}\label{eq:H}
  H(\rho):=\frac{\Lambda(\rho)}{\rho^{2}}
   =\sum_{k\ge1}\frac{\rho^{2k-2}}{k(2k-1)}\quad(0<\rho<1),
  \qquad H(0):=1 .
\end{equation}
The series converges for $|\rho|<1$, so $H$ is continuous, strictly increasing and
$\ge1$ on $[0,1)$, and $\Lambda(\rho)=\rho^{2}H(\rho)$ there.

\begin{theorem}[Large-argument asymptotics]\label{thm:asym}
Let $x_{j},y_{j}>0$, $v_{j}=\tfrac{x_j+y_j}{2}\to\infty$, $u_{j}=\tfrac{x_j-y_j}{2}$, and
put $\rho_{j}:=\frac{|u_{j}|}{1+v_{j}}$. Suppose
\[
  \rho_{j}\;\longrightarrow\;\rho_{0}\in[0,1).
\]
Then, with $H$ as in \eqref{eq:H},
\[
  \log\Gs(x_{j},y_{j})
   =\frac{u_{j}^{2}}{1+v_{j}}\,H(\rho_{j})\big(1+o(1)\big)
   =(1+v_{j})\,\Lambda(\rho_{j})\big(1+o(1)\big)
  \qquad (j\to\infty),
\]
and the analogous statement holds for $\log\Gg$ with $1+v_j$ replaced by $v_j$
throughout, i.e.\ under the hypothesis $|u_j|/v_j\to\rho_0\in[0,1)$ one has
$\log\Gg=\frac{u_j^{2}}{v_j}H\big(\frac{|u_j|}{v_j}\big)(1+o(1))$. (The sign of $u_j$ is
immaterial, since $\Lambda$ and $H$ are even.)
\end{theorem}

The expansion extends verbatim to any number of variables.

\begin{theorem}[Multivariate expansion]\label{thm:multi}
Let $n\ge2$, $x_{1},\dots,x_{n}>0$, $\bar x=\frac1n\sum_{i=1}^{n}x_{i}$ and
$m_{j}=\sum_{i=1}^{n}(x_{i}-\bar x)^{j}$ for $j\ge2$. If
$\max_{i}|x_{i}-\bar x|<\bar x$, then
\begin{equation}\label{eq:multi}
  \log\frac{\prod_{i=1}^{n}\Gamma(x_{i})}{\Gamma(\bar x)^{n}}
   =\sum_{j=2}^{\infty}\frac{(-1)^{j}}{j}\,\zeta(j,\bar x)\,m_{j},
\end{equation}
the series converging absolutely. In particular
$\prod_{i=1}^{n}\Gamma(x_{i})\ge\Gamma(\bar x)^{n}$, with equality iff $x_{1}=\dots=x_{n}$.
\end{theorem}

Finally we address Open Problem 3. We first note that $t$ is uniquely determined, then we
control its distance to the midpoint by two independent arguments, and we determine the
sharp constant in the resulting bound.

\begin{theorem}[Localization of $t$]\label{thm:loc}
Let $x\ne y$ be positive, $v=\tfrac{x+y}{2}$, $s=\sqrt{xy}$, $c=\tfrac{|x-y|}{2}$,
$A=1+v$ and $\rho=c/A$. Then:
\begin{enumerate}
\item[\rm(i)] there is a unique $t=t(x,y)\in(s,v)$ satisfying \eqref{eq:Wt};
\item[\rm(ii)] with $\rho=\rho(x,y)$ one has
  \begin{equation}\label{eq:locbnd}
    \frac{c^{2}\,\zeta(4,A)}{4\,\zeta(3,A)+6v\,\zeta(4,A)}
    \;\le\; v-t \;\le\;
    \frac{c^{2}\,\zeta(4,A)}{4\,\zeta(3,A)}\Big(1+\frac{2\rho^{2}}{3(1-\rho^{2})}\Big);
  \end{equation}
\item[\rm(iii)] as $c\to0^{+}$ with $v$ fixed,
  \begin{equation}\label{eq:locasym}
    v-t=\frac{c^{2}\zeta(4,A)}{4\zeta(3,A)}
     +\frac{c^{4}}{2\zeta(3,A)}\Big[\frac{\zeta(6,A)}{3}
       -\frac{3\,\zeta(4,A)^{3}}{16\,\zeta(3,A)^{2}}\Big]+O(c^{6});
  \end{equation}
\item[\rm(iv)] $t>\frac{1}{2}\big(\sqrt{xy}+\frac{x+y}{2}\big)$ whenever at least one of
  the following two conditions holds:
  \begin{align}
    &\text{(A)}\qquad \frac{v+s}{2(1+v)}\Big(1+\frac{2\rho^{2}}{3(1-\rho^{2})}\Big)<1,
      \label{eq:condA}\\[1mm]
    &\text{(B)}\qquad \mathcal{E}(x,y):=
      \int_{v+\frac12}^{\infty}\!\!K-\frac{1}{A-w}-\frac{1}{2(A-w)^{2}}\;<\;0,
      \label{eq:condB}
  \end{align}
  where $w=\frac{v-s}{2}$ and $K(\xi)=\frac{1}{c^{2}}\log\frac{1}{1-c^{2}/\xi^{2}}$
  for $\xi>c$.
\end{enumerate}
\end{theorem}

Condition (A) is a Hurwitz-zeta comparison; condition (B) is a convexity bound, and
$\int_{B}^{\infty}K$ has the closed form
\begin{equation}\label{eq:intK}
  \int_{B}^{\infty}K=\frac{1}{c^{2}}\Big[-B\log\frac{B^{2}}{B^{2}-c^{2}}
     +c\log\frac{B+c}{B-c}\Big],\qquad B>c .
\end{equation}
The two conditions are complementary --- (A) is effective for moderate $\rho$, (B) as
$\rho\to1$ --- and numerically their union covers the whole parameter range
(Section \ref{sec:num}). This leads us to sharpen the second half of Open Problem 3.

\begin{theorem}\label{conj:sharp}
For all $x\ne y>0$,
\begin{equation}\label{eq:conj}
  \frac{t(x,y)-\sqrt{xy}}{\frac{x+y}{2}-\sqrt{xy}}\;>\;\frac{2}{3},
\end{equation}
equivalently $t(x,y)>\frac13\big(x+y+\sqrt{xy}\big)$. The constant $\frac23$ is best
possible and the bound is not attained: the ratio in \eqref{eq:conj} tends to $\frac23$
along $v\to\infty$, $|x-y|/(x+y)\to0$.
\end{theorem}

\begin{proof}
Fix $v>0$, let $0<c<v$, $s=\sqrt{v^{2}-c^{2}}$, $A=1+v$, $\Delta=\frac{v-s}{3}$ and
$\Phi(c)=\log\Gs(x,y)/c^{2}=\sum_{k\ge1}\frac{c^{2k-2}}{k}\zeta(2k,A)$. Since
$\zeta(2,1+t)=\Phi(c)$ and $\zeta(2,\cdot)$ is strictly decreasing,
\[
  \theta>\tfrac23\iff t>\tfrac{2v+s}{3}=v-\Delta
  \iff\zeta(2,A-\Delta)>\Phi(c)\iff R(c)>0 ,
\]
and $R(0)=0$ by continuity. We prove $R'(c)>0$ on $(0,v)$. Because
$\Delta'=\frac{c}{3s}$ and $\partial_{\Delta}\zeta(2,A-\Delta)=2\zeta(3,A-\Delta)$,
\begin{equation}\label{eq:Rp}
  R'(c)=\frac{2c}{3s}\zeta(3,A-\Delta)-\frac{2}{c}E(c),\qquad
  E(c)=\sum_{m\ge1}\frac{m}{m+1}c^{2m}\zeta(2m+2,A),
\end{equation}
so it suffices to prove
\begin{equation}\label{eq:target}
  E(c)<\frac{c^{2}}{3s}\zeta(3,A-\Delta).
\end{equation}

\emph{Step 1 (a global midpoint bound for $E$).} Put
$K_c(u)=\frac{1}{c^{2}}\log\frac{1}{1-c^{2}/u^{2}}$ and
$Q_c(u)=\frac{1}{u^{2}-c^{2}}-K_c(u)$. From
$K_c(u)=\sum_{j\ge0}\frac{c^{2j}}{j+1}u^{-2j-2}$ and
$\frac{1}{u^{2}-c^{2}}=\sum_{j\ge0}c^{2j}u^{-2j-2}$,
\[
  Q_c(u)=\sum_{m\ge1}\frac{m}{m+1}c^{2m}u^{-2m-2}>0,\qquad E(c)=\sum_{n\ge0}Q_c(A+n).
\]
All coefficients being positive, $Q_c$ is positive, decreasing and convex on $u>c$; the
midpoint rule for a convex function, summed over $n\ge0$, gives with $B=A-\frac12>c$
\begin{equation}\label{eq:mid}
  E(c)<\int_{B}^{\infty}Q_c(u)\,du ,
\end{equation}
with no restriction on the size of $c^{2}/A$. Term-by-term integration with
$y=c^{2}/B^{2}\in(0,1)$ gives
\begin{equation}\label{eq:F}
  \int_{B}^{\infty}Q_c(u)\,du=\frac{c^{2}}{6B^{3}}F(y),\qquad
  F(y)=6\sum_{m\ge1}\frac{m}{(m+1)(2m+1)}y^{m-1}.
\end{equation}

\emph{Step 2 (a sharp one-variable inequality).} Let
$S(y)=\sum_{m\ge0}\frac{y^{m}}{(m+1)(2m+1)}$, so $F=6S'$, and let
$h(y)=\frac{1-\sqrt{1-y}}{2+\sqrt{1-y}}$, for which a direct differentiation gives
$6h'(y)=\frac{9}{\sqrt{1-y}\,(2+\sqrt{1-y})^{2}}$. The identity
$h=\frac{3(1-\sqrt{1-y})-y}{3+y}$, together with
$1-\sqrt{1-y}=\sum_{m\ge1}c_my^{m}$,
$c_m=\frac{\binom{2m}{m}}{4^{m}(2m-1)}$, $\frac{c_m}{c_{m-1}}=\frac{2m-3}{2m}$, shows
that $h=\sum_{m\ge1}a_my^{m}$ with $a_1=\frac16$ and
$a_m=c_m-\frac13a_{m-1}$ for $m\ge2$. Writing $S=1+\sum_{m\ge1}b_my^{m}$,
$b_m=\frac{1}{(m+1)(2m+1)}$, we have $a_1=b_1=\frac16$ and, for $r_m=a_m/c_m$,
$r_m=1-\frac{2m}{3(2m-3)}r_{m-1}$ with $r_2=\frac59$, $r_3=\frac{17}{27}$; a two-line
induction gives $\frac{3(m-1)}{4m}\le r_m<\frac34$ for every $m\ge2$. Hence
$a_2-b_2=\frac{1}{360}$, $a_3-b_3=\frac{11}{3024}$, and for $m\ge4$
$a_m>c_m-\frac14c_{m-1}=\frac{3(m-2)}{2(2m-3)}c_m$, while
$Q_m:=\frac{3(m-2)c_m}{2(2m-3)b_m}$ satisfies $Q_4=\frac{135}{128}>1$ and
$\frac{Q_{m+1}}{Q_m}=\frac{(m-1)(m+2)(2m-3)(2m+3)}{2(m-2)(m+1)^{2}(2m+1)}>1$ because
$2m^{3}-5m^{2}+5m+22>0$. Thus $a_m>b_m$ for all $m\ge2$, so $h'>S'$ and
\begin{equation}\label{eq:sharp1}
  F(y)<\frac{9}{\sqrt{1-y}\,(2+\sqrt{1-y})^{2}},\qquad 0<y<1 .
\end{equation}

\emph{Step 3 (the finite-$v$ correction only helps).} Let $z=A-\Delta$,
$e=\frac{1}{2B}\in(0,1)$ and $q=\frac{s}{B}$, so that
$q=\sqrt{(1-e)^{2}-y}$ and, exactly, $z=\frac{B(2+q+4e)}{3}$. A direct computation gives
\[
  \frac{B^{3}(z+1)}{sz^{3}}=L_e(y):=\frac{9\,(2+q+10e)}{q\,(2+q+4e)^{3}},
\]
and $L_0(y)$ is precisely the right-hand side of \eqref{eq:sharp1}. Since
$q'(e)=-\frac{1-e}{q}$,
\[
  \frac{d}{de}\log L_e=-\frac{N(e,q)}{q^{2}(2+q+4e)(2+q+10e)},
\]
where
$N(e,q)=40e^{3}+40e^{2}q-12e^{2}+83eq^{2}-32eq-24e+2q^{3}+q^{2}-8q-4$. For fixed
$e\in(0,1)$ the function $q\mapsto N(e,q)$ is strictly convex
($\partial_{q}^{2}N=166e+12q+2>0$), and on $0<q<1-e$ its endpoint values are
$N(e,0)=4(e-1)(2e+1)(5e+1)<0$ and
$N(e,1-e)=9(e-1)(9e^{2}-2e+1)<0$. A convex function lies below the maximum of its
endpoint values, so $N(e,q)<0$ on $0<q<1-e$; hence $e\mapsto L_e(y)$ is strictly
increasing and
\begin{equation}\label{eq:sharp2}
  \frac{9}{\sqrt{1-y}\,(2+\sqrt{1-y})^{2}}<L_e(y)=\frac{B^{3}(z+1)}{sz^{3}} .
\end{equation}

\emph{Step 4 (conclusion).} Combining \eqref{eq:mid}--\eqref{eq:sharp2},
\[
  E(c)<\frac{c^{2}}{6B^{3}}F(y)<\frac{c^{2}}{6B^{3}}\cdot\frac{B^{3}(z+1)}{sz^{3}}
        =\frac{c^{2}(z+1)}{6sz^{3}} .
\]
Finally, the trapezoid rule applied to the convex decreasing function
$t\mapsto(z+t)^{-3}$ gives
$\zeta(3,z)\ge\int_{0}^{\infty}(z+t)^{-3}dt+\frac{1}{2z^{3}}=\frac{z+1}{2z^{3}}$, so
$E(c)<\frac{c^{2}}{3s}\zeta(3,z)$, which is \eqref{eq:target}. Therefore $R'(c)>0$ on
$(0,v)$; with $R(0)=0$ this gives $R(c)>0$, and by \eqref{eq:Rp} and the first display
$\theta>\frac23$, as claimed. \qed
\end{proof}

\begin{conjecture}\label{conj:ab}
Conditions (A) and (B) of Theorem \ref{thm:loc} are satisfied for every $x\ne y>0$.
\end{conjecture}

\begin{theorem}[An explicit certified range]\label{thm:partial}
Let $x\ne y>0$ satisfy
\begin{equation}\label{eq:partialcond}
  (x-y)^{2}\;\le\;\frac{2+x+y}{3},
\end{equation}
equivalently $c^{2}\le\frac A6$, where $c=\frac{|x-y|}{2}$ and $A=1+v$. Then
$\theta(x,y)>\frac23$. In particular, for every fixed $d\ne0$ inequality \eqref{eq:conj}
holds for all $x,y>0$ with $|x-y|=d$ and $x+y$ sufficiently large: this is precisely the
range in which the constant $\frac23$ is sharp.
\end{theorem}

\begin{remark}[A product formula]\label{rem:product}
The Weierstrass product for $\Gamma$ gives a representation of $\Gs$ that makes the
identity $\Phi=\sum_{m\ge0}K(A+m)$ transparent. Indeed
$\Gamma(z)=\frac{e^{-\gamma z}}{z}\prod_{n\ge1}\frac{e^{z/n}}{1+z/n}$ yields, with
$v=\frac{x+y}{2}$ and $c=\frac{|x-y|}{2}$,
\[
  \frac{\Gamma(v+c)\Gamma(v-c)}{\Gamma(v)^{2}}
   =\frac{v^{2}}{v^{2}-c^{2}}\prod_{n\ge1}\frac{(n+v)^{2}}{(n+v)^{2}-c^{2}} ,
\]
the factors $e^{\pm z/n}$ cancelling. Since
$\Gs(x,y)=\frac{4xy}{(x+y)^{2}}\Gg(x,y)=\frac{v^{2}-c^{2}}{v^{2}}\cdot
 \frac{\Gamma(v+c)\Gamma(v-c)}{\Gamma(v)^{2}}$, the factors $v^{2}-c^{2}$ cancel as well
and
\begin{equation}\label{eq:product}
  \Gs(x,y)=\prod_{n\ge1}\frac{(n+v)^{2}}{(n+v)^{2}-c^{2}}
          =\prod_{m\ge0}\frac{(A+m)^{2}}{(A+m)^{2}-c^{2}} ,
\end{equation}
Taking logarithms, and writing $A=1+v$, recovers $\log\Gs=\sum_{m\ge0}\log\big(1-\frac{c^{2}}{(A+m)^{2}}\big)^{-1}
=c^{2}\sum_{m\ge0}K(A+m)=\sum_{k\ge1}\frac{c^{2k}}{k}\zeta(2k,A)$, i.e.\
Theorem~\ref{thm:main} for $\Gs$; each factor of \eqref{eq:product} is positive because
$c<v<A$. In particular $\Gs$ is a product of terms $>1$, which is the Jensen--Gurland
inequality in multiplicative form. The individual summand in $R(c)$ is positive once
$u>\frac{c^{2}+\Delta^{2}}{2\Delta}$: from $-\log(1-y)\le\frac{y}{1-y}$ with
$y=\frac{c^{2}}{u^{2}}$ it is enough that $u^{2}-c^{2}>(u-\Delta)^{2}$, i.e.\
$u>\frac{c^{2}+\Delta^{2}}{2\Delta}$. So only the
$\max\big\{0,\lfloor\frac{c^{2}+\Delta^{2}}{2\Delta}-A\rfloor+1\big\}$ smallest indices can be
exceptional, and
numerically their total deficit is exceeded by the total surplus of the others by a factor
that falls to $1.0016$ at $v=10^{4}$, $c=0.1v$. This is the quantitative reason why
estimates coarser than about $10^{-3}$ relative accuracy cannot decide
Theorem~\ref{conj:sharp}.
\end{remark}

\begin{remark}[A monotone formulation]\label{rem:monotone}
Theorem \ref{conj:sharp} would also follow from the combination of the following two
statements about the
single quantity $\Theta(c):=\dfrac{v-t(c)}{v-s(c)}$, $0<c<v$, where $t(c)$ is the point
with $\zeta(2,1+t)=\log\Gs/c^{2}$ and $s=\sqrt{v^{2}-c^{2}}$:
\begin{itemize}
\item[(m1)] $\Theta$ is strictly decreasing on $(0,v)$;
\item[(m2)] $\lim_{c\to0^{+}}\Theta(c)\le\frac13$.
\end{itemize}
Indeed (m1) and (m2) give $\Theta(c)<\frac13$ for every $c>0$, which is
$v-t<\frac{v-s}{3}$, i.e.\ $\theta>\frac23$. Statement (m2) is a local computation: as
$c\to0$ one has $v-t\sim\frac{c^{2}\zeta(4,A)}{4\zeta(3,A)}$ by \eqref{eq:locasym} and
$v-s\sim\frac{c^{2}}{2v}$, so
\[
  \lim_{c\to0^{+}}\Theta(c)=\frac{v\,\zeta(4,A)}{2\,\zeta(3,A)}
   =\frac13\Big(1-\frac{1}{2A}+O(A^{-2})\Big)\qquad(A\to\infty),
\]
and this limit is $<\frac13$ for every $A>1$: by \eqref{eq:ratio},
$\frac{\zeta(4,A)}{\zeta(3,A)}\le\frac{2A^{3}}{3(A-\frac12)^{3}(A+1)}$, so
$\lim_{c\to0^{+}}\Theta(c)<\frac13$ is equivalent to $(A-1)A^{3}<(A-\frac12)^{3}(A+1)$,
i.e.\ to $4A^{3}-6A^{2}+5A-1>0$, which holds for $A\ge1$. Statement
(m1) is what remains: numerically it holds for every tested $v$ ($1.5$, $2$, $5$, $20$,
$100$, $10^{3}$, $10^{4}$), $\Theta$ decreasing from $\frac13-\frac{1}{6A}$ at $c\to0$ to
values as low as $0.183$ at $c\to v$ for $v=\frac32$. Equivalently, by
$\zeta(2,A-\delta)=\Phi$ and $\delta'=E/(c\,\zeta(3,A-\delta))>0$, while
$\Delta'=c/(3s)$, so (m1) reads
\begin{equation}\label{eq:monotone}
  \delta(c)\,c^{2}\,\zeta(3,A-\delta(c))\;>\;s(c)\,(v-s(c))\,E(c),
  \qquad E(c)=\sum_{m\ge1}\frac{m}{m+1}c^{2m}\zeta(2m+2,A).
\end{equation}
\end{remark}

\begin{remark}[A wider certified range]\label{rem:wider}
The proof of Theorem \ref{thm:partial} can be repeated with a sharper split of the
logarithm. Since $-\log(1-y)=y+\frac{y^{2}}{2}+\sum_{k\ge3}\frac{y^{k}}{k}
\le y+\frac{y^{2}}{2}+\frac{y^{3}}{3(1-y)}$ for $0<y<1$, and since
$(u-\Delta)^{2}\le u^{2}$, one has
\[
  c^{2}R(c)\;\ge\;c^{2}\big[2\Delta\zeta(3,A)-\Delta^{2}\zeta(4,A)\big]
   -\frac{c^{4}\zeta(4,A)}{2(1-\rho^{2})} ,
\]
so $R(c)>0$ certainly holds as soon as
\begin{equation}\label{eq:wider}
  2\Delta\,\zeta(3,A)\;>\;\Delta^{2}\zeta(4,A)+\frac{c^{2}\zeta(4,A)}{2(1-\rho^{2})},
  \qquad \Delta=\frac{v-s}{3},\ \rho=\frac{c}{A}.
\end{equation}
Condition \eqref{eq:wider} is strictly weaker than the hypothesis of
Theorem \ref{thm:partial}: it holds throughout
$c^{2}\lesssim0.62\,A$ (numerically $c^{2}/A=0.8145$, $0.6661$, $0.6297$, $0.6216$,
$0.6208$ for $v=5,20,100,1000,10^{4}$), whereas $c^{2}\le\frac A6\approx0.167A$ was
all that the earlier estimates gave. Beyond $c^{2}\approx0.62A$ the right-hand side of
\eqref{eq:wider} overtakes the left and this particular chain stops discriminating.
\end{remark}

\begin{remark}[Scaling limit]\label{rem:scaling}
Fix $\rho=\frac{c}{A}$ and let $A\to\infty$. Since $\zeta(2m+2,A)\sim\frac{A^{-(2m+1)}}{2m+1}$
and $v-s=A(1-\sqrt{1-\rho^{2}})(1+o(1))$, the comparison $\zeta(2,A-\Delta)>\Phi$ behind
Theorem \ref{conj:sharp} becomes, at leading order in $A^{-1}$,
\begin{equation}\label{eq:scaling}
  \frac{t}{1-t}\;>\;\Sigma_{0}(\rho)-1,
  \qquad t=\frac{1-\sqrt{1-\rho^{2}}}{3},
  \qquad \Sigma_{0}(\rho)=\sum_{m\ge0}\frac{\rho^{2m}}{(m+1)(2m+1)} ,
\end{equation}
and $\Sigma_{0}$ has the closed form
$\Sigma_{0}(\rho)=2\frac{\arsh\rho}{\rho}+\frac{\log(1-\rho^{2})}{\rho^{2}}$. The difference
of the two sides of \eqref{eq:scaling},
$D(\rho)=\frac{t}{1-t}-\Sigma_{0}(\rho)+1$, is positive on $(0,1)$, and this is again a
consequence of Step 2 of the proof of Theorem \ref{conj:sharp}. Indeed, with $y=\rho^{2}$,
\[
  \frac{t}{1-t}=\frac{1-\sqrt{1-y}}{2+\sqrt{1-y}}=h(y),
  \qquad \Sigma_{0}(\rho)=S(y),
  \qquad\text{so } D=h-S+1 .
\]
Step 2 writes $h=\sum_{m\ge1}a_{m}y^{m}$ and $S=1+\sum_{m\ge1}b_{m}y^{m}$ with
$a_{1}=b_{1}=\frac16$ and $a_{m}>b_{m}$ for every $m\ge2$; hence
$D=\sum_{m\ge2}(a_{m}-b_{m})y^{m}$ has positive coefficients, which proves $D>0$ and at the
same time identifies the expansion
$D(\rho)=\frac{\rho^{4}}{360}+\frac{11\rho^{6}}{3024}+O(\rho^{8})$, since $a_2-b_2=\frac1{360}$
and $a_3-b_3=\frac{11}{3024}$ are the coefficients computed there (with a positive
$O(\rho^{8})$ tail). Numerically $D(\rho)\approx\rho^{4}/360$ down to $\rho=5\cdot10^{-4}$
(where $D=1.74\cdot10^{-16}$) and $D(0.99)=5.10\cdot10^{-2}$. Thus the constant $\frac23$ is
decided by an $O(\rho^{4})$ cancellation: this is why no first-order estimate settles
Theorem \ref{conj:sharp}.
\end{remark}

\begin{remark}\label{rem:equiv}
Theorem \ref{conj:sharp} can be restated as a comparison of $\zeta(2,\cdot)$ with its
own triangular average. With $A=1+v$, $c=|x-y|/2$, $K$ as in \eqref{eq:intK} and
$\Phi=\log\Gs(x,y)/c^{2}$, the identity
$\sum_{n\ge0}K(A+n)=\Phi$ of the proof of Theorem \ref{thm:loc}(iv) combines with
$\psi(A+t)-\psi(A-t)=\int_{-t}^{t}\psi'(A+z)\,dz$ to give
\[
  \Phi=\frac{1}{c^{2}}\int_{0}^{c}\big[\psi(A+t)-\psi(A-t)\big]dt
     =\int_{-c}^{c}\psi'(A+z)\,\frac{c-|z|}{c^{2}}\,dz .
\]
Since $\zeta(2,1+t)=\psi'(1+t)$, the inequality of Theorem \ref{conj:sharp} is therefore
equivalent to
\[
  \int_{-c}^{c}\Big[\psi'\!\big(A-\Delta\big)-\psi'(A+z)\Big]\frac{c-|z|}{c^{2}}\,dz>0,
  \qquad \Delta=\frac{v-s}{3}=\frac{c^{2}}{3(v+s)},
\]
where $s=\sqrt{xy}$: it compares the value of $\psi'=\zeta(2,\cdot)$ at the point
$A-\Delta$ with its average over $[-c,c]$ against the triangular weight $(c-|z|)/c^{2}$,
whose variance is $c^{2}/6$. Expanding both sides about $A$, the terms of order $c^{0}$ and
$c^{1}$ cancel (the weight is even), and the first surviving comparison is
\[
  \frac{\zeta(4,A)}{\zeta(3,A)}<\frac{4}{3(v+s)} ,
\]
which holds for all $x,y>0$, with relative margin $\frac{1}{A}$ because
$\zeta(4,A)/\zeta(3,A)=\frac{2}{3A}\big(1+O(A^{-1})\big)$ and $s\le v$. The inequality is
delicate because the decision rests on the next terms, of relative size $\rho^{2}$: this is
also why the constant $\frac23$ is sharp but not attained.
\end{remark}

\begin{remark}
The limiting statement in Theorem \ref{conj:sharp} follows from Theorem
\ref{thm:loc}(iii). Take $x_{j}=v_{j}(1+\eps_{j})$, $y_{j}=v_{j}(1-\eps_{j})$ with
$v_{j}\to\infty$, $\eps_{j}\to0$ and $v_{j}\eps_{j}\to0$. Then $c_{j}=v_{j}\eps_{j}\to0$,
so \eqref{eq:locasym} applies and, since $\zeta(4,A)/\zeta(3,A)\sim2/(3A)$ as
$A\to\infty$,
\[
  v_{j}-t_{j}\sim\frac{v_{j}^{2}\eps_{j}^{2}}{4}\cdot\frac{2}{3(1+v_{j})}
  \sim\frac{v_{j}\eps_{j}^{2}}{6},
  \qquad
  v_{j}-s_{j}=\frac{v_{j}\eps_{j}^{2}}{1+\sqrt{1-\eps_{j}^{2}}}\sim\frac{v_{j}\eps_{j}^{2}}{2}.
\]
Hence $(v_{j}-t_{j})/(v_{j}-s_{j})\to\frac13$ and $(t_{j}-s_{j})/(v_{j}-s_{j})\to\frac23$. Thus
$\theta=\frac23+\frac{1}{6v}+O(v^{-2})$; the penultimate column of Table \ref{tab:inf} is
such a path and confirms this to nine significant digits. It also shows that a fixed
$\eps$ does \emph{not} suffice: there $\theta\to\frac23+g(\eps)>\frac23$.
\end{remark}
